\section{Introduction}

Quantum key distribution (QKD) moves the security argument from computational hardness to physics. Its information-theoretic guarantee holds only inside a device model; once the hardware departs from that model, an eavesdropper can learn the key while every monitored statistic looks normal \cite{scarani2009,xu2020}. Detector blinding is the standard example. Bright illumination turns single-photon avalanche detectors into classical threshold detectors, and a faked-state eavesdropper then dictates which detector clicks \cite{makarov2005,lydersen2010}; a full-field implementation against a running link left the quantum bit error rate (QBER) unchanged and gave the attacker the entire key \cite{gerhardt2011}.

Operators therefore watch telemetry. A deployed decoy-state BB84 link \cite{bennett2014,hwang2003,lo2005,wang2005} reports, per monitoring window, detection counts and error counts for every intensity and basis, and long field trials already feed some of these signals back into polarisation and phase stabilisation \cite{dixon2010,dynes2012,stucki2011}. A recent strand of work trains machine-learning (ML) classifiers on such telemetry and reports near-perfect detection of intercept-resend, photon-number-splitting (PNS), Trojan-horse and blinding attacks \cite{mao2020,liao2022,alkuwari2026,isha2026,mohamed2026}. These results come from simulators. In most of them the benign class is stationary or randomised sample by sample, the classes are balanced, and the attacker is not calibrated against the statistics the defender monitors; the adversarial formulation of \cite{mohamed2026}, with domain randomisation and block-level decisions, is the main exception and the closest work (Section~\ref{sec:background}).

Two questions follow that the cited studies do not settle. The first concerns the attacker: which attacks are visible in window counts at all, given that a capable eavesdropper can shape the statistics she produces? The second concerns the defender: how many false alarms does a telemetry detector raise when the fibre drifts for a day, as installed fibre does \cite{dixon2010,dynes2012}, rather than sitting at its calibration point? A detector that alarms every few minutes on a healthy link will be switched off; a detector that misses the one attack that breaks the security proof adds nothing that the decoy-state parameter estimation does not already provide.

This paper studies both questions for three-intensity decoy-state BB84 with the finite-key analysis of Lim \etal \cite{lim2014}, restricted to \emph{count telemetry}: per-window pulse, click, sifted-click and error counts with Alice's intensity-monitor readings. The boundary of what such telemetry can show is elementary. An attacker who reproduces the honest photon-number yields and error probabilities leaves the law of the counts unchanged (Proposition~\ref{prop:indist}), and the perfect eavesdropper of \cite{gerhardt2011}, demonstrated on an entanglement-based system without decoy states, is its analogue. The contributions lie in what happens inside that boundary.
\begin{itemize}
\item A simulation benchmark with Bob's basis choice and sifting, per-basis misalignment, afterpulses, basis-dependent efficiencies, slow intensity drift and an intensity monitor, run over 24-hour streams at five drift scales. It covers three in-model attacks and ten faked-state attacks, from a naive one to one that emulates the honest channel exactly; one needs only an attenuated threshold detector, one estimates the channel from public counts, two ignore part of the drift, one stays just inside physically plausible parameters, and one is intermittent. All design choices were frozen before the reported runs, which use seeds not seen during development. Every attacked window is compared with the same window of an honest copy, and every post-processing block is audited against a ground-truth count of the key bits Eve knows.
\item A comparison of thirteen detectors, including an adaptive-reference baseline, that separates two design choices: what a detector is calibrated on, and whether its features are invariant to drift. Features built on a per-window refit of the honest channel, the decoy-consistency statistic (DCS) and detectors trained on its residuals, keep the false-alarm rate near target at every drift scale with drift-free calibration only, under two stated preconditions on the link. Level-based detectors need drift-calibrated thresholds and then miss the target in both directions.
\item A map of what count telemetry sees along the attacker ladder. Naive and intermittent blinding are caught at once. Error-free and background-free emulation is caught only by sequential or drift-calibrated detectors. Emulation with plausible parameters, emulation from public estimates and exact emulation pass every detector; threshold-detector emulation is separated only by a post-hoc block-level statistic.
\item Two side results with practical weight. A quantitative illustration of a known requirement: with nominal intensities, slow intensity drift of a few per cent makes the single-photon bound of \cite{lim2014} fail in a sizeable share of honest blocks. And physical plausibility floors on the fitted parameters, an obvious remedy for the blind spots of the refit, fail on ordinary operational changes and are evaded by an attacker who stays above them.
\end{itemize}

The scope is deliberately narrow. All data are simulated, drift amplitudes are assumptions, attacks are pulse-wise independent, and the claims are about what count telemetry can establish. Richer observables such as timing and coincidence statistics \cite{dasilva2012,gras2021} and device-level countermeasures \cite{yuan2010,lim2015,lo2012} are outside the detectors studied here and are discussed in Section~\ref{sec:discussion}.

Section~\ref{sec:background} reviews decoy-state QKD and prior work on quantum hacking and its detection. Section~\ref{sec:model} fixes the channel and threat models, and Section~\ref{sec:limits} states the indistinguishability boundary. Section~\ref{sec:dcs} defines the DCS. Sections~\ref{sec:setup} and~\ref{sec:results} give the experimental design and the results, and Section~\ref{sec:discussion} discusses limitations.
